\documentclass{article}
\usepackage{graphicx}
\usepackage{amsmath}
\usepackage{amssymb}
\usepackage{amsfonts}
\usepackage{hyperref}
\usepackage{exercises}
\usepackage{indentfirst}
\usepackage{circledsteps}
\usepackage{amsthm}
\usepackage{csquotes}
\usepackage{yhmath}
\usepackage{comment}
\usepackage{placeins}
\usepackage{geometry}
\usepackage{hyperref}
\usepackage{tikz}
\newcommand*{\oldneg}{\mathord{\sim}}
\newcommand{\Z}{\mathbb{Z}}
\newcommand{\N}{\mathbb{N}}
\let\oldemptyset\emptyset

\theoremstyle{definition}
\newtheorem{definition}{Definition}
\newtheorem{theorem}{Theorem}[section]
\newtheorem{corollary}{Corollary}[theorem]
\newtheorem{lemma}[theorem]{Lemma}

\setlength{\parindent}{.5in}
\geometry{margin=0.5in}
\setlength{\parskip}{12pt}
\title{Chapter 5 Notes}
\begin{document}
\date{}
\maketitle

\section*{Section 1}

\subsection*{Notes}
After Corollary 1.9 it reads: \textquote{We thus see that the function $\sigma$ is strictly progressing.} Certainly, this was shown for 
\textit{ordinals}, but we have not shown this for any sets in general. In fact, $\sigma$ being strictly progressing for all sets 
is equivalent to $x\in x$ not holding for any set $x$. For, on the one hand, if $\sigma$ is strictly progressing for all sets, 
suppose $x\in x$ for some set $x$. Now, $x^+=x\cup\{x\}$. Also, $x\subset x^+$ (by assumption). But consider any element $y$
of $x^+$. Either $y\in x$ or $y=x$. In either case, because we assumed $x\in x$, $y\in x$. Thus $x^+\subseteq x$, contrary to the 
assumption that $x\subset x^+$. Hence $x\ne x$ for all sets $x$. On the other hand, suppose that $x\ne x$ for all sets 
$x$. Then it must be that $x\subset x^+$. For obviously $x\subseteq x^+$, and if $x=x^+$, then, because $x\in x^+$, 
$x\in x$, a contradiction. Thus $\sigma$ would be strictly progressing.

Hence $\sigma$ being strictly progressing for all sets is equivalent to $x\in x$ holding for no sets, which we have not shown.

In the proof of Theorem 1.17, it is said that no two distinct elements (of $On$) can have the same successor. To prove this, we 
first prove the lemma that (for any two ordinals), if $\beta<\alpha$, then $\beta^+\ne\alpha^+$. For suppose 
$\beta<\alpha$ and $\beta^+=\alpha^+$. Thus $\alpha\cup\{\alpha\}=\beta\cup\{\beta\}$. Thus either $\alpha=\beta$ or 
$\alpha\in\beta$. The first possibility is ruled out by the assumption that $\beta<\alpha$. Also, if $\alpha\in\beta$, the, 
by Theorem 1.12, $\alpha<\beta$, contrary to $\beta<\alpha$. Thus $\alpha\notin\beta^+$, a contradiction. Hence 
$\beta^+\ne\alpha^+$. 

Now suppose $\alpha^+=\beta^+$. Thus $\alpha\cup\{\alpha\}=\beta\cup\{\beta\}$. Thus, if $x\in\alpha$, then 
either $x\in\beta$ or $x=\beta$. If $x=\beta$, then, because $\alpha$ and $\beta$ are ordinals (and by Theorem 1.12), 
$\beta<\alpha$. But, by the above lemma, this means that $\beta^+\ne\alpha^+$, a contradiction. Hence $x\ne\beta$, meaning 
$x\in\beta$ (and thus $\alpha\subseteq\beta$). Similar reasoning shows that $\beta\subseteq\alpha$, meaning 
$\alpha=\beta$, therefore completing the proof.

In Theorem 1.19, it was remakred that $\alpha<\beta$ implies that $\alpha^+\le\beta$. For suppose that 
$\alpha<\beta$. Then, by Theorem 1.12, $\alpha\in\beta$. Hence $\alpha\in\beta$ and for all $x\in\alpha$, 
$x\in\beta$, meaning $\alpha^+\le\beta$.

In Theorem 1.20, it is remarked that $\beta<\alpha^+$ implies that $\beta\le\alpha$. Here is why that is: For suppose that 
$\beta<\alpha^+$. Thus every element of $\beta$ is either an element of $\alpha$, or is equal to $\alpha$. Suppose 
some element $x$ of $\beta$ was equal to $\alpha$. Because every ordinal is transitive, that would mean, for every 
$y\in\alpha$, $y\in\beta$. Hence $\beta=\alpha\cup\{\alpha\}$, contrary to the fact that $\beta<\alpha^+$. Thus all elements of 
$\beta$ are elements of $\alpha$ (and none are equal to $\alpha$) meaning $\beta\le\alpha$.

\section*{Section 2}

\subsection*{Notes}
Note that Theorem 2.3 presupposes that $x$ is a set, as this means that there is at least one ordinal outside of 
$x$ (because $On$ is not a set).

\subsection*{Exercise 2.1}
This is an easy consequence of Theorem 2.2. For, by Theorem 2.2, every transitive subset of $On$ is an element of $On$, and 
thus $On$ is $T$-closed. Hence, if $x$ belongs to all $T$-closed classes, then $x\in On$, meaning $x$ is an ordinal.

\subsection*{Exercise 2.2}
This is another easy consequence of Theorem 2.2. For consider any non-empty set $A$ of ordinals. 
Let $B=\cap A$. Obviously, $B$ is a set of ordinals. It suffices to show that $B$ is transitive (and hence is an 
ordinal by Theorem 2.2). Consider any $x\in B$. Thus, for every $y\in A$, $x\in y$. But, because $y$ is an 
ordinal (and thus it is transitive), all elements $z$ of $x$ are in $y$. Thus all elements $z$ of $x$ are 
in all elements of $A$, meaning all elements $z$ of $x$ are in $\cap A$. Thus $B$ is transitive, 
meaning it is an ordinal.

\subsection*{Exercise 2.3}
(a): On the one hand, we show that every ordinal $x$ (which is of course a set) 
can be well ordered in such a way that each element $y$ of 
$x$ is the set of all elements of $x$ that precede $y$ in the ordering. We show that every ordinal 
can be well ordered under $\subseteq$ so that the required conditions are satisfied. We prove this 
using superinduction. First, recall that $On$ is well ordered under $\subseteq$.\newline
(1): Clearly $\emptyset$ is well ordered under $\subseteq$ such that the required condition is 
satisfied (in fact, $\subseteq\restriction\emptyset=\emptyset$).\newline
(2): Now suppose $\alpha$ is well ordered such that the required condition is satisfied. Now consider 
$\alpha^+$. It has already been shown that $\alpha^+$ is a transitive set of ordinals. Thus, because 
it is a subclass of $On$, it is well ordered under $\subseteq$ (because $On$ is well ordered under 
$\subseteq$). We need to show that $\subseteq\restriction\alpha^+$ satisfies the required condition. 
If $x\in\alpha^+$, either $x\in\alpha$ or $x=\alpha$ (because $\alpha^+=\alpha\cup\{\alpha\}$). 
If $x\in\alpha$, then $x$ is the set of all lesser elements in $\alpha$ by assumption. But it has 
already been shown that $x<\alpha$ (because $x\in\alpha$ iff $x<\alpha$), meaning 
$x$ is the set of all elements of $\alpha^+$ less than $x$. If $x=\alpha$, then obviously 
$x$ is the set of all lesser elements of $\alpha^+$. Hence, for every $x\in\alpha^+$, $x$ is the set of all 
lesser elements of $\alpha^+$.\newline
(3): Lastly, suppose that some chain $C$ of ordinals is well ordered under $\subseteq$ and satisfies
the required conditions. Thus, by assumption, every element $x\in C$ is the set of all lesser elements 
of $C$. Consider any element $y\in\cup C$. Then it is an element of some $x\in C$, and thus, by assumption, 
it is also an element of $C$. Hence, by assumption, $y$ is the set of all lesser elements of $C$. Thus every 
element in $\cup C$ is an element of $C$, meaning $\cup C\subseteq C$ and thus $\cup C$ is well ordered 
under $\subseteq$ (as $C$ was assumed to be well ordered under $\subseteq$). Now we show that for any $y\in\cup C$ (and 
recall that we have also shown that $y\in C$), the 
set of all elements in $\cup C$ less than $y$ is equal to the set of all elements in $C$ less than $y$. For suppose 
$x$ is an element in $\cup C$ less than $y$. Then, as we have shown, $x$ is also an element in $C$. Thus $x\in C$, meaning $x$ is 
a lesser element of $y$ that is in $C$. 
On the other hand, if $x$ is an element in $C$ less than $y$, then, because $y\in C$, $x\in y$, meaning $x\in\cup C$. Thus we 
have shown that, for every $y\in\cup C$, it is the set of all lesser elements in $C$, and also the set of all lesser 
elements in $\cup C$. Thus $\subseteq\restriction\cup C$ satisfies the required condition.

Hence, by superinduction, it follows that every ordinal is well ordered and satisfies the given condition under 
$\subseteq$. 

(b): On the other hand, we show that every set $x$ that can be well ordered under $\le$ in such a way that each element $y$ of $x$ 
is the set of all elements of $x$ that precede $y$ in the ordering, is an ordinal. It can be shown by transfinite induction 
on $x$ that every element $y$ of $x$ is an ordinal. For suppose that for all elements $z$ less than $y$, $z$ is an ordinal. 
Then, by definition of $x$, $y$ is the set of all lesser elements of $x$. Thus $y$ is a set of ordinals. And 
$y$ is transitive (because if $z\in y$ and $v\in z$, then $z\in x$, and thus $v<z$ and $z<y$, meaning 
$v<y$ (because $\le$ is transitive) and thus $v\in y$). Thus $y$ is a transitive set of ordinals, meaning it is an ordinal.
Thus every element of $x$ is an ordinal. Obviously $x$ is transitive (because, by definition of $x$, if 
$y\in x$ and $z\in y$, then $z<y$ and $z\in x$). Thus $x$ is a transitive set of ordinals, meaning it is an ordinal. 
This completes the proof.

\section*{Section 3}

\subsection*{Notes}
Let $\omega\times 1=\omega$, and, for all $x>1$, $\omega\times x+0=\omega\times x$. As the book notes, $\omega\times x+y$ being a set 
for every $x>0$ and every $y$ can be proven using the axiom of substitution. So, for this section, assume the axiom of substitution, so that we know that 
$\omega\times x+y$ is a set for every $x>0$ and every $y$. In this section, we prove some additional results about these sets.
First, we prove that for every 
$x>0$ and every $y$, $\omega\times x+y$ is transitive, and $\omega\times x+y\notin\omega\times x+y$. 

an ordinal. Recall that
\[ \omega\times n=0,1,2,...,\omega,\omega+1,\omega+2,...,\omega\times 2,\omega\times 2+1,\omega\times 2+2,...,\omega\times (n-1),\omega\times (n-1)+1,\omega\times (n-1)+2,... \]
\begin{lemma} $\omega\times n+m=\omega\times n\cup\{\omega\times n,...,\omega\times n+(m-1)\}$ for all $m>0$.
\end{lemma}
\begin{proof}
The base case, when $m=1$, follows by definition, as $\omega\times n+1=\omega\times n\cup\{\omega\times n\}$. Now supose the result 
holds for some $m>0$. Then
\[ \omega\times n+(m+1)=\omega\times n+m\cup\{\omega\times n+m\} \]
\[ =\cup\{\omega\times n,\{\omega\times n,...,\omega\times n+(m-1)\}\}\cup\{\omega\times n+m\} \]
\[ =\omega\times n\cup\{\omega\times n,...,\omega\times n+(m-1),\omega\times n+m\} \]
therefore completing the induction. 
\end{proof}

\begin{lemma} If for some $x$, if for $m<x$, then $\omega\times m+n\in\omega\times x$ for every $n$,
then for every $y$, if $m<x$ or $m=x$ and $n<y$, then $\omega\times m+n\in\omega\times x+y$.
\end{lemma}
\begin{proof}
Suppose the hypothesis. We prove that the result holds by induction on $y$. The base case holds
by assumption. Now suppose that the result holds for some $y$. Then
\[ \textrm{(1) }\omega\times x+(y+1)=\omega\times x+y\cup\{\omega\times x,...,\omega\times x+y\} \] by 
Lemma 0.1. By assumption, all sets of the form $\omega\times m+n$ where $m<x$ are elements of
$\omega\times x$, and hence also elements of $\omega\times x+(y+1)$ (by (1)). Also, all sets of the form
$\omega\times m+n$ where $m=x$ and $n<y$ are elements of $\omega\times x+(y+1)$ by $(1)$. Hence the result holds
for $y+1$, meaning the induction is complete.
\end{proof}

Note that this above lemma implies that for every $x>0$ and every $y$, if $m<x$ or $m=x$ and $n<y$, then
$\omega\times m+n\in\omega\times x+y$, because, by definition, for every $m<x$, then
$\omega\times m+n\subset\omega\times x$, meaning the hypothesis of the lemma is satisfied for every $x$. And of course the converse holds
as well: every element of $\omega\times x+y$ is either a natural number, or a set of the form
$\omega\times m+n$ where either $m<x$ or $m=x$ and $n<y$.

Now, note that, because $\omega$ is a transitive set of ordinals, it is also an ordinal.
\begin{theorem} For every $x>0$ and every $y$, $\omega\times x+y$ is transitive.
\end{theorem}
\begin{proof}
By Lemma 0.2, we know that that, for every $x>1$ and every $y$, $\omega\times x+y$ is just the set of all natural numbers, 
as well as the set of all $\omega\times m+n$, where either $0<m<x$, or $m=x$ and $n<y$. Now consider any element $e$ of an element of 
$\omega\times x+y$. Thus it is either an element of a natural number (in which case it is also a natural number, and hence also 
a member of $\omega\times x+y$), or an element of the form $\omega\times m+n$ where either $0<m<x$ or $m=x$ and $n<y$. In the latter case, 
if it is not a natural number, it must be an element of the form $\omega\times v+w$ where either $0<v<m$ or $v=m$ and $w<n$. But that means 
it would be of the form $\omega\times v+w$ where either $0<v<x$ or $v=x$ and $w<y$, meaning $e$ is an element of $\omega\times x+y$. Hence 
$\omega\times x+y$ is transitive.
\end{proof}

\begin{lemma} If $\omega\times x\notin\omega\times x$, then $\omega\times x+y\notin\omega\times x+y$ for every $y$.
\end{lemma}
\begin{proof}
Suppose $\omega\times x\notin\omega\times x$. We prove by induction that $\omega\times x+y\notin\omega\times x+y$ for every $y$.
The base case, when $y=0$, is given by the assumption. Now suppose the result holds for some $y$. 
Then $\omega\times x+(y+1)=\omega\times x+y\cup\{\omega\times x+y\}$. If $\omega\times x+(y+1)\in\omega\times x+(y+1)$, then 
either $\omega\times x+(y+1)=\omega\times x+y$, or $\omega\times x+(y+1)\in\omega\times x+y$. In the first case, we would 
have $\omega\times x+y\in\omega\times x+y$, a contradiction. In the second case, because, by Theorem 0.3, 
$\omega\times x+y$ is transitive,
$\omega\times x+y\in\omega\times x+y$, a contradiction. Thus
$\omega\times x+(y+1)\notin\omega\times x+(y+1)$, therefore completing the proof. 
\end{proof}

\begin{theorem} For every $x>0$ and every $y$, $\omega\times x+y\notin\omega\times x+y$.
\end{theorem}
\begin{proof}
We prove this by induction on $x$. The base case is when $x=1$, that is, when 
$\omega\times x+y=\omega$. But $\omega\notin\omega$. Thus, by Lemma 0.4, 
$\omega+y\notin\omega+y$ for all $y$. Thus the base case holds. Now supose the result holds for all numbers less 
than $x$, for $x>1$. Then $\omega\times x$ is the set of all natural numbers, together with 
every $\omega\times m+n$ for some $m<x$. So suppose $\omega\times x\in\omega\times x$. Obviously it cannot be that 
$\omega\times x$ was a natural number, meaning $\omega\times x=\omega\times m+n$ for some $m<x$. But then, because 
$\omega\times m+n$ is transitive, $\omega\times m+n\in\omega\times m+n$, contradicting the induction assumption. Hence 
$\omega\times x\notin\omega\times x$, and consequently, $\omega\times x+y\notin\omega\times x+y$ for every $y$, as a consequence 
of Lemma 0.4. 
\end{proof}

Assuming that, for every $x>0$ and every $y$, $\omega\times x+y$ is a set, we now prove that it is also an ordinal.

\begin{lemma} If $\omega\times x$ is an ordinal, then $\omega\times x+y$ is an ordinal for every $y$.
\end{lemma}
\begin{proof}
Suppose the hypothesis. Then we prove, by induction on $y$, that $\omega\times x+y$ is an ordinal for every $y$. The 
base case, when $y=0$, holds by assumption. Now suppose the result holds all numbers less than $y$, for $y>0$. By 
Lemma 0.1, 
\[ \omega\times x+y=\omega\times x\cup\{\omega\times x,\omega,...,\omega\times x+(y-1)\} \] Now, because 
$\omega\times x$ was assumed to be an ordinal, it follows that it is the set of all lesser ordinals, meaning all of its elements are also ordinals. 
Also, by the induction assumption, all elements of $\{\omega\times x,\omega,...,\omega\times x+(y-1)\}$ are ordinals, meaning 
all elements of $\omega\times x+y$ are ordinals. By Theorem 0.5, $\omega\times x+y$ is transitive. Thus it is a transitive set of 
ordinals, meaning it is also an ordinal.
\end{proof}

\begin{theorem} For every $x>0$ and every $y$, $\omega\times x+y$ is an ordinal.
\end{theorem}
\begin{proof}
This can be proven by induction on $x$, proving that the property: $\omega\times x+y$ holds for every $y$, holds for every $x>0$. 
The base is when $x=1$, that is, when 
$\omega\times x=\omega$. But $\omega$ is an ordinal (because it is a transitive set of ordinals). Thus, by Lemma 0.6,
$\omega+y$ is an ordinal for every $y$. Now suppose the result holds for all numbers less than $x$, for $x>1$. Then 
$\omega\times x$ is just the set of all natural numbers (which are all ordinals), and all sets of the form $\omega\times m+n$ where $m<x$, 
which are ordinals by the induction assumption. Thus $\omega\times x$ is a transitive set of ordinals, meaning it is also an ordinal, 
therefore completing the induction.
\end{proof}

Lastly, we show that all sets of the form $\omega\times n$ are limit ordinals, whereas all sets of the form 
$\omega\times n+m$ where $m>0$ are successor ordinals.

\begin{theorem} (1) All ordinals of the form $\omega\times n$ are limit ordinals. (2) All ordinals of the form
$\omega\times n+m$ for $m>0$ are successor ordinals, in particular, they are successors of the ordinal 
$\omega\times n+(m-1)$.
\end{theorem}
\begin{proof}
(1) Consider some ordinal of the form $\omega\times n$. Clearly it is not the least ordinal. Now, suppose it was a successor ordinal. Hence, 
for some $\alpha$, $\alpha<\omega\times n$ and for no $\beta$ is it the case that $\alpha<\beta<\omega\times n$. But, because 
$\omega\times n$ is an ordinal, $\alpha\in\omega\times n$. Thus it is either a natural number $a$, or of the form 
$\omega\times x+y$ where $x<n$. In the first case, then $a<a^+<\omega\times x$, a contradiction to the assumption that 
$\omega\times n$ was the successor of $\alpha$. In the second case, then 
$\omega\times x+y<\omega\times x+(y+1)<\omega\times n$, another contradiction. Thus $\omega\times n$ cannot be a successor ordinal, 
and is consequently a limit ordinal.

(2) Now consider some ordinal of the form $\omega\times n+m$ for $m>0$. We show that 
$\omega\times n+m$ is the successor of $\omega\times n+(m-1)$. Clearly $\omega\times n+(m-1)<\omega\times n+m$. But 
suppose there was some $\alpha$ such that $\omega\times n+(m-1)<\alpha<\omega\times n+m$. Now, 
$\omega\times n+m$ has exactly one more element than $\omega\times n+(m-1)$, namely $\omega\times n+(m-1)$ itself. Thus, 
it must be that $\alpha=\omega\times n+m$, a contradiction. Hence there can be no $\alpha$ such that 
$\omega\times n+(m-1)<\alpha<\omega\times n+m$, meaning $\omega\times n+m$ is the successor ordinal of 
$\omega\times n+(m-1)$.
\end{proof}

\end{document}
